% Source: https://artofproblemsolving.com/community/c3h2793844_2018_geometry_bee__guts_rounds__sets_14
% Stable identity: aops-topic-2793844-q08
Write the side length as $L$. The two centre coordinates each differ by at most $L-3r$, so their distance is at most $\sqrt2(L-3r)$. Nonoverlap requires distance at least 3 r. Hence $r\le\dfrac{\sqrt2L}{3(1+\sqrt2)}=\dfrac{2}{3}$. Equality is attained by placing the centres at $(r,r)$ and $(L-2r,L-2r)$, where the disks are tangent. Thus $\boxed{\dfrac{2}{3}}$.
\begin{asy}
usepackage("amsmath");
size(7cm);
real L=2+sqrt(2),r=2.0/3;pair A=(r,r),B=(L-2*r,L-2*r);draw((0,0)--(L,0)--(L,L)--(0,L)--cycle);draw(circle(A,r));draw(circle(B,2*r));draw(A--B,dashed);label("$r$",A,S);label("$2r$",B,N);
\end{asy}
